% ACL 2024 Style Paper: Rhetorical Framing Auditor
% Final experimental results and analysis

\documentclass[11pt]{article}
\usepackage{acl}
\usepackage{times}
\usepackage{latexsym}
\usepackage{amsmath}
\usepackage{amssymb}
\usepackage{booktabs}
\usepackage{multirow}
\usepackage{graphicx}
\usepackage{xcolor}
\usepackage{subcaption}

\title{Fact Omission vs. Rhetorical Framing: \\
What Predicts Media Bias?}

\author{Anonymous Submission}

\begin{document}
\maketitle

\begin{abstract}
We investigate the relative contributions of \textbf{fact omission} and \textbf{rhetorical framing} in predicting media bias across left, center, and right news sources. Using a dataset of 1,748 news triplets covering the same stories, we apply Rhetorical Structure Theory (RST) analysis to extract structural features and develop the Discourse Framing Index (DFI) to quantify how different sources position shared facts. Our experiments reveal a striking finding: \textbf{fact omission accounts for approximately 72\% of predictive power}, while RST-based structural framing contributes only 28\%. When we restrict analysis to facts reported by all three bias sides (eliminating the omission signal), classification accuracy drops from 89.77\% to 61.46\%---barely above random chance. These results suggest that media bias manifests primarily through \textit{what} facts sources choose to report, rather than \textit{how} they structurally position shared facts.
\end{abstract}

\section{Introduction}

Media bias detection has traditionally focused on lexical and semantic features---word choice, sentiment, and framing language. However, bias can also manifest through structural choices: which facts to include or omit, and where to position them within the discourse hierarchy. This work investigates a fundamental question: \textbf{Is media bias better predicted by fact omission (selection bias) or by the rhetorical positioning of shared facts (framing bias)?}

We hypothesize that these two mechanisms contribute differently to detectable bias:
\begin{itemize}
    \item \textbf{Tier 1: Selection Bias (Omission)} --- Sources choose which facts to report, creating coverage asymmetries
    \item \textbf{Tier 2: Framing Bias (RST Structure)} --- For facts that are shared across sources, the rhetorical positioning differs
\end{itemize}

Our experimental journey began with the hypothesis that RST-based structural framing would be a strong predictor of bias. Through systematic experimentation, we discovered that this hypothesis was largely incorrect: omission dominates the predictive signal.

\section{Related Work}

Prior work on media bias detection has examined lexical features \cite{}, sentiment analysis \cite{}, and more recently, discourse-level features. RST \cite{mann1988rhetorical} provides a framework for analyzing how text segments relate hierarchically. Our work builds on the AllSides Media Bias dataset, which provides triplets of articles covering the same story from left, center, and right sources.

\section{Methodology}

\subsection{Dataset and RST Parsing}

We use 1,748 news article triplets from the AllSides dataset, each containing articles from left-leaning, center, and right-leaning sources covering the same news event. Each article is parsed using a state-of-the-art RST parser (UniRST) to extract Elementary Discourse Units (EDUs) and their hierarchical relationships.

After RST parsing, we obtain:
\begin{itemize}
    \item 630,516 total EDUs across all articles
    \item 548,960 EDUs retained after quality filtering (removing boilerplate, URLs, etc.)
    \item Depth and nuclearity information for each EDU
\end{itemize}

\subsection{Fact Clustering}

To identify shared facts across articles, we cluster semantically similar EDUs using:
\begin{enumerate}
    \item \textbf{SBERT Encoding}: EDUs are encoded using \texttt{all-mpnet-base-v2}
    \item \textbf{Agglomerative Clustering}: Cosine distance with average linkage (threshold 0.35)
    \item \textbf{Cross-Encoder Validation}: Pairs validated using \texttt{ms-marco-MiniLM-L-6-v2} (threshold 0.7)
\end{enumerate}

This produces 16,170 total clusters across all triplets. Critically, we observe the following distribution:
\begin{itemize}
    \item \textbf{All 3 sides present}: 3,037 clusters (18.8\%)
    \item \textbf{Left + Center only}: 7,057 clusters (43.6\%)
    \item \textbf{Center + Right only}: 6,076 clusters (37.6\%)
    \item \textbf{Left + Right only}: 0 clusters (0.0\%)
\end{itemize}

The absence of Left+Right clusters (without Center) is by design---we require a Center EDU for baseline comparison in DFI calculation.

\subsection{The Bipartite Approach}

Our finalized feature extraction method uses a \textbf{bipartite decomposition} approach that creates pairings between bias sides within each cluster.

\subsubsection{Matching Process}

For each cluster containing EDUs from multiple bias sides:
\begin{enumerate}
    \item Extract EDUs grouped by bias: $E_L$ (left), $E_C$ (center), $E_R$ (right)
    \item Perform greedy 1-to-1-to-1 matching to create triplets $(e_l, e_c, e_r)$
    \item Matching criterion: minimize depth difference across the triplet
    \item Remaining unmatched EDUs represent coverage asymmetries
\end{enumerate}

\subsubsection{Feature Vector Construction}

For each matched triplet and unmatched EDU, we compute:

\textbf{Coverage Features:}
\begin{equation}
\delta_{cov}^{L} = \begin{cases}
0 & \text{if matched} \\
+1 & \text{if left-only (left has, center doesn't)} \\
-1 & \text{if center-only (center has, left doesn't)}
\end{cases}
\end{equation}

\textbf{Structural Features:}
\begin{equation}
\delta_{str}^{L} = W(e_l) - W(e_c)
\end{equation}
where $W(e) = \frac{1}{1 + \log(1 + depth(e))}$ is the prominence weight.

The same formulation applies symmetrically for right vs. center comparisons.

\subsubsection{Handling Missing Matches}

When a cluster lacks EDUs from one or more bias sides:
\begin{itemize}
    \item \textbf{No left EDU}: Left features receive $\delta_{cov}^L = -1$ (center has fact, left doesn't)
    \item \textbf{No right EDU}: Right features receive $\delta_{cov}^R = -1$
    \item \textbf{No center EDU}: Cluster is excluded (center required as baseline)
\end{itemize}

This design means that \textbf{coverage features encode the omission signal}---when one side omits a fact that another side reports.

\subsection{Cluster Ordering Strategies}

The DFI vector concatenates features from all clusters. The ordering of clusters affects the feature vector structure. We experimented with multiple ordering strategies:

\begin{enumerate}
    \item \textbf{Deterministic (by cluster ID)}: Clusters ordered by their integer ID from the clustering algorithm
    \item \textbf{By min depth}: Clusters ordered by the minimum normalized EDU depth (shallowest = most prominent first)
    \item \textbf{By max prominence}: Clusters ordered by maximum prominence score
    \item \textbf{By cluster size}: Clusters ordered by number of EDUs
\end{enumerate}

Note: The ``deterministic'' ordering is not random---it reflects the order in which the agglomerative clustering algorithm formed clusters, which implicitly captures semantic similarity patterns.

\subsection{Classification Models}

We experimented with two classifier families:

\textbf{Support Vector Machine (SVM):} Initial experiments used SVM with RBF kernel (C=10, $\gamma$=0.1). However, SVM requires fixed-length input vectors, necessitating padding/truncation of variable-length DFI vectors.

\textbf{Random Forest:} We transitioned to Random Forest (200 trees, max depth 15) because:
\begin{enumerate}
    \item Better handling of variable-importance patterns across positions
    \item No sensitivity to padding artifacts
    \item Built-in feature importance for interpretability
    \item Superior performance in our experiments
\end{enumerate}

\section{Experiments and Results}

\subsection{Experiment 1: Full Dataset Classification}

Using the complete dataset with all cluster configurations:

\begin{table}[h]
\centering
\begin{tabular}{lcc}
\toprule
\textbf{Method} & \textbf{Test Acc} & \textbf{F1} \\
\midrule
Bipartite Coverage + RF & \textbf{89.77\%} & 0.897 \\
Padded DFI + RF & 87.78\% & 0.877 \\
Padded DFI + SVM & 85.23\% & 0.852 \\
\bottomrule
\end{tabular}
\caption{Full dataset classification results}
\label{tab:full-results}
\end{table}

\subsection{Experiment 2: RST-Only Features}

To isolate the RST structural contribution, we trained models using only structural features (no coverage):

\begin{table}[h]
\centering
\begin{tabular}{lcc}
\toprule
\textbf{Feature Set} & \textbf{Test Acc} & \textbf{F1} \\
\midrule
Full RST + Logistic Regression & \textbf{61.16\%} & 0.602 \\
Depth features only & 58.43\% & 0.571 \\
Role features only & 54.21\% & 0.532 \\
Satellite count only & 56.89\% & 0.558 \\
\bottomrule
\end{tabular}
\caption{RST-only classification (no coverage signal)}
\label{tab:rst-only}
\end{table}

\subsection{Experiment 3: Quantifying Omission vs. Framing}

Comparing full features vs. RST-only:

\begin{equation}
\text{RST Contribution} = \frac{61.16\% - 50\%}{89.77\% - 50\%} \approx 28\%
\end{equation}

\begin{equation}
\text{Omission Contribution} = 100\% - 28\% = 72\%
\end{equation}

This calculation assumes 50\% as the random baseline for binary classification (left-vs-center or right-vs-center).

\subsection{Experiment 4: Cluster Ordering Ablation}

Testing whether semantic ordering of clusters improves classification:

\begin{table}[h]
\centering
\begin{tabular}{lcc}
\toprule
\textbf{Ordering Strategy} & \textbf{Test Acc} & \textbf{$\Delta$ vs Baseline} \\
\midrule
Deterministic (cluster ID) & \textbf{89.77\%} & --- \\
Min depth (shallowest first) & 88.92\% & -0.85\% \\
Max prominence & 88.14\% & -1.63\% \\
Cluster size & 87.56\% & -2.21\% \\
\bottomrule
\end{tabular}
\caption{Cluster ordering comparison (bipartite coverage features)}
\label{tab:ordering}
\end{table}

Surprisingly, the deterministic ordering outperforms all ``semantically meaningful'' orderings. This suggests the implicit ordering from agglomerative clustering captures useful structure.

\subsection{Experiment 5: Pure 3-Way Cluster Analysis}

To definitively test whether RST framing alone can predict bias (without any omission signal), we created a dataset containing \textbf{only clusters where all 3 bias sides are present}.

\subsubsection{3-Way Clustering Methodology}

We implemented an anchor-based approach:
\begin{enumerate}
    \item Use center EDUs as anchors
    \item For each center EDU, find best matching left AND right EDU (cosine threshold 0.65, cross-encoder threshold 0.5)
    \item Only retain clusters where all 3 matches are found
\end{enumerate}

\subsubsection{3-Way Dataset Statistics}

\begin{itemize}
    \item Total triplets: 1,748
    \item Triplets with $\geq 1$ 3-way cluster: 1,321 (75.6\%)
    \item \textbf{Triplets with zero 3-way clusters: 427 (24.4\%)}
    \item Total 3-way clusters: 5,216
    \item \textbf{Average 3-way clusters per triplet: 3.95}
\end{itemize}

The low average (3.95 clusters/triplet) and 427 triplets with zero 3-way clusters demonstrates that \textbf{left and right sources rarely report the exact same facts}---they systematically omit different information.

\subsubsection{3-Way Classification Results}

\begin{table}[h]
\centering
\begin{tabular}{lcc}
\toprule
\textbf{Method} & \textbf{Test Acc} & \textbf{vs Full Dataset} \\
\midrule
\multicolumn{3}{l}{\textit{Deterministic ordering:}} \\
Bipartite structural & \textbf{61.46\%} & -28.31\% \\
Bipartite combined & 61.46\% & -28.31\% \\
\midrule
\multicolumn{3}{l}{\textit{Min-depth ordering:}} \\
Bipartite structural & 57.29\% & -32.48\% \\
Bipartite combined & 58.33\% & -31.44\% \\
\midrule
Coverage only & \textbf{50.00\%} & -39.77\% \\
\bottomrule
\end{tabular}
\caption{3-way cluster classification (no omission signal)}
\label{tab:3way-results}
\end{table}

\textbf{Critical observation:} Coverage-only features achieve exactly 50\% (random chance) because all coverage deltas are zero---every cluster has exactly one EDU from each bias side.

\subsection{Experiment 6: Structural Ablation}

Using the standard dataset, we ablated RST structural components:

\begin{table}[h]
\centering
\begin{tabular}{lccc}
\toprule
\textbf{Experiment} & \textbf{$\alpha$} & \textbf{$\gamma$} & \textbf{Test Acc} \\
\midrule
Baseline (full structure) & 0.8 & 0.5 & 89.77\% \\
Without satellite effect & 0.8 & 1.0 & 89.12\% \\
Without depth effect & 1.0 & 0.5 & 88.45\% \\
Without both & 1.0 & 1.0 & 87.23\% \\
\bottomrule
\end{tabular}
\caption{Structural ablation results}
\label{tab:structural-ablation}
\end{table}

Removing structural weighting reduces accuracy by ~2.5\%, confirming RST features provide modest but measurable signal.

\section{Discussion}

\subsection{The Two-Tier Model of Media Bias}

Our experiments support a two-tier model:

\textbf{Tier 1: Selection Bias (72\%)}
\begin{itemize}
    \item Sources choose which facts to report
    \item Left and right sources omit different facts
    \item Only 18.8\% of clusters have all 3 bias sides
    \item Average of only 3.95 shared facts per triplet
\end{itemize}

\textbf{Tier 2: Framing Bias (28\%)}
\begin{itemize}
    \item For shared facts, sources position them differently
    \item RST depth and nuclearity provide weak signal
    \item Best pure-framing accuracy: 61.46\% (barely above chance)
\end{itemize}

\subsection{Why Deterministic Ordering Outperforms}

The deterministic (cluster ID) ordering unexpectedly outperforms semantically-motivated orderings. We hypothesize:
\begin{enumerate}
    \item Agglomerative clustering forms clusters in order of semantic density
    \item Earlier clusters (lower IDs) may represent more ``core'' facts
    \item This implicit ordering captures journalistic salience better than RST depth
\end{enumerate}

\subsection{Why We Chose Random Forest Over SVM}

Initial experiments used SVM, but we transitioned to Random Forest for several reasons:
\begin{enumerate}
    \item \textbf{Variable-length handling}: DFI vectors vary in length across triplets; RF handles this more gracefully than SVM's required padding
    \item \textbf{Feature importance}: RF provides interpretable feature importance scores
    \item \textbf{Performance}: RF achieved 89.77\% vs. SVM's 85.23\%
    \item \textbf{Robustness}: RF is less sensitive to hyperparameter choices
\end{enumerate}

\subsection{Implications for Media Bias Detection}

Our findings suggest that:
\begin{enumerate}
    \item \textbf{Fact coverage analysis} may be more valuable than framing analysis for bias detection
    \item \textbf{Simple binary features} (fact present/absent) capture most of the signal
    \item \textbf{RST-based methods} provide incremental improvements but are not the primary signal
\end{enumerate}

\section{Conclusion}

We set out to investigate whether RST-based rhetorical framing could predict media bias. Through systematic experimentation, we discovered that \textbf{fact omission is the dominant predictor}, accounting for approximately 72\% of classification accuracy. When restricted to facts shared across all bias sides, classification drops to 61.46\%---barely above random chance.

This finding has important implications: media bias manifests primarily through \textit{editorial selection}---what facts to include or omit---rather than through the \textit{structural positioning} of shared facts. Future work on media bias detection should prioritize coverage analysis alongside traditional framing features.

\section{Reproducibility}

All code and experimental configurations are available at [anonymous repository]. Key parameters:
\begin{itemize}
    \item SBERT model: \texttt{all-mpnet-base-v2}
    \item Cross-encoder: \texttt{ms-marco-MiniLM-L-6-v2}
    \item Clustering threshold: 0.35 (cosine distance)
    \item Random Forest: 200 trees, max depth 15
    \item Train/val/test split: 80\%/10\%/10\%
\end{itemize}

\bibliographystyle{acl_natbib}
\bibliography{references}

\end{document}
